\pdfminorversion=7
\documentclass[12pt,a4paper]{article}

\input{"./../references/preamble.tex"}
\input{"./../references/macros.tex"}

\usepackage{titling}
\usepackage[compact]{titlesec}
\usepackage[a4paper,top=0.8in,left=0.8in,right=0.8in,bottom=0.8in]{geometry}
\usepackage{cite}
\usepackage{wrapfig}
\usepackage{subcaption}
\usepackage[multiple]{footmisc}
\usepackage{lmodern}

\usepackage[numbers]{natbib}
\setlength{\bibsep}{0pt plus 0.3ex}

\usepackage{fancyhdr}  % use this package to get a 2 line header
\pagestyle{fancy}
\fancyhf{}
\fancyhead[L]{\it Google India PhD Fellowship 2026\/}
\fancyhead[R]{{\bf Srinjoy Sarkar}}
\fancyfoot{}
\fancyfoot[R]{\thepage}

\renewcommand{\headrulewidth}{1.0pt}
\renewcommand{\footrulewidth}{0.4pt}
\setlength{\headheight}{18pt}

\usepackage{hyperref}
\hypersetup{
  breaklinks, % Allow links to break across lines
  hidelinks, % Set to true to hide colored boxes and use plain text links
  urlcolor=black, % Color for URLs
  % Border color if boxes are used (only if colorlinks=false)
  linkbordercolor=white,
  % Links will be colored, the default color is red. False will use boxed links.
  colorlinks=true,
  % Internal links, those generated by cross-referenced elements, are
  % displayed in blue.
  linkcolor=blue,
  % Links to local files will be shown in magenta color (see linking local files)
  filecolor=magenta,
  hyperfootnotes=false}

\newcommand{\htuvm}{\textsf{HT-UVM}\xspace}
\newcommand{\Htuvm}{\textsf{HT-UVM}\xspace}

\newcommand{\htovs}{\textsf{HT-OVS}\xspace}
\newcommand{\Htovs}{\textsf{HT-OVS}\xspace}

\newcommand{\skipuvm}{\textsf{SL-UVM}\xspace}
\newcommand{\Skipuvm}{\textsf{SL-UVM}\xspace}

\newcommand{\skipovs}{\textsf{SL-OVS}\xspace}
\newcommand{\Skipovs}{\textsf{SL-OVS}\xspace}

\newcommand{\heterosl}{\textsf{SL-OVS}\xspace}

\setlength{\droptitle}{-2cm}
\title{Exploring Efficient GPU Data Structures for Oversubscription}
\author{Srinjoy Sarkar \\
    Department of Computer Science Engineering, \\
Indian Institute of Technology Kanpur, India,  \\
\texttt{srinjoys23@cse.iitk.ac.in}}

\date{} %leave this blank

\begin{document}
\maketitle
\thispagestyle{fancy}

\section*{Introduction}

The high compute capability and memory bandwidth of GPUs, coupled with energy
efficiency, have spurred their usage across a diverse set of application domains
including scientific computing simulations, big data and large-scale graph
analytics, and computational biology~\cite{nvidia-applications, tensorflow-2015, metacache-icpp-2021}.
Furthermore, improved compute capabilities have led to a steady rise in the
working set sizes of GPU-accelerated applications, which now often far exceed
the GPU global memory capacity~\cite{dynamap,uvm-forest-isca25}. Hence, the
traditional copy-then-execute (CTE) programming model is giving way to the newer
unified virtual memory (UVM) paradigm, where the CPU virtual address space is
also visible to GPU kernels. GPU vendors have introduced memory virtualization
support through features such as unified virtual addressing, demand paging, and
prefetching. UVM allows the host CPU and discrete GPU devices to share a unified
virtual address space, enabling the GPU runtime to page memory in and out as
needed~\cite{hum-ppopp-2020,dynamap}. Demand paging in UVM-capable GPUs
facilitates transparent memory oversubscription, improving programmer
productivity.
% The cost of servicing GPU page faults can be reduced by
% prefetching a chunk of nearby
% pages~\cite{cppe-ipdps-2020,uvmsmart-isca-2019,uvm-forest-isca25,hpe-tcad-2020}.

Writing correct but scalable multithreaded GPU kernels is a difficult art~\cite{massively-parallel-processors-v4,kernelbench-arxiv}. In addition to the algorithm and the compute capabilities of the GPU device, the scalability of the application depends on the building blocks of the kernel, \eg, GPU data structures. For example, the performance of the Pinterest application improved substantially by switching to GPU-accelerated hash tables~\cite{pinterest-perf-boost}. Therefore, the design of sophisticated GPU data structures such as lists and trees is an active area of research~\cite{gpu-gems2}. Much recent work has focused on designing efficient data structures (\eg, hash tables, skip lists, B-trees, and membership query filters) for use in device code~\cite{slabhash-ipdps-2020,dachash-sbacpad-2021,warpcore-hipc-2020,warpdrive,parallel-hashmap-nvidia,survey-hashing-tpds-2020,bght,dycuckoo,gph-pamd-2025,gfsl-pact-2017,cmedia-skiplist-ipdpsw-2019,skiplist-survey-2025,multiversion-gpu-btree,gpu-filters-ppopp23,harmonia-gpu,gpu-btree,gpu-gelhash,survey-hashing-tpds-2020,lockfree-ds-icpads-2012,quotient-filter,stdgpu,cucollections-nvidia,gpu-cds}.

\paragraph*{The Problem}

The steady increase in the working set size and irregular access patterns of GPU applications require more efficient memory management so that the GPU data structures efficiently scale with the input size. However, existing studies primarily focus on the traditional CTE model (\ie, in-memory data structures) and do not address how to design efficient data structures that oversubscribe memory, and developers fall back to UVM. While UVM simplifies programming, its performance remains a significant concern due to frequent far faults\footnote{Far faults, which occur when a requested page is absent in GPU global memory, require multiple PCIe roundtrips, and are resolved after interacting with the CPU page tables. On \NVIDIA GPUs, a far fault from a CUDA thread stalls the progress of the warp, possibly affecting the overall GPU throughput.},
high address translation overheads, and high volume of page replacement~\cite{suv-micro-2024,in-depth-analysis-uvm-sc21,fine-grain-quantitative-uvm-taco24,batch-aware-asplos-2020,uvmsmart-isca-2019,cppe-ipdps-2020,irregular-page-walks-isca-2018,heterogeneous-uvm,uvm-forest-isca25,dynamap,uvmbench-arxiv}. For instance, Lin et al.~\cite{oversubscription-dl-training} find UVM-based implementations competitive with the CTE model for workloads that fit in GPU memory, but about $7\times$ slower with oversubscription. The performance degradation with UVM is \emph{more} prominent in applications with irregular access patterns and pointer chasing~\cite{fine-grain-quantitative-uvm-taco24,demystifying-uvm-cost-ipdps2021,in-depth-analysis-uvm-sc21}.

Data structures like hash tables, skip lists, B-trees, and approximate membership query (AMQ) filters are characterized by input-dependent access locality~\cite{art-of-multiprocessor-programming-v2}.
% As a result,
Lack of locality in the input operation strings and suboptimal implementation
choices can lead to irregular access patterns and poor data locality. When
memory oversubscription is coupled with these performance pathologies, these
data structures can exhibit very poor scalability. In this proposal, we
establish that UVM-based implementations of popular random-access data
structures with GPU memory oversubscription do not scale well with input size
even when these implementations are meticulously optimized using data usage
hints and prefetches based on memory/translation
usage~\cite{nvidia-prefetching}. Our experiments show that the run times of
these implementations increase exponentially with a linear increase in the input
sizes indicating a clear \emph{performance bottleneck}. We focus on alleviating
this performance pathology by proposing efficient designs of popular concurrent data
structures for GPUs (\eg, hash table, skip list, AMQ filter, and B-trees).

\section*{Brief Literature Review}

Hash table is a popular concurrent data structures for storing key-value
pairs~\cite{art-of-multiprocessor-programming}. Hash tables are widely used in
various scientific applications like metagenomic
sequencing~\cite{metacache-icpp-2021}, photodna~\cite{photodna-ares-2023}, and
database indexing workloads that are accelerated using the GPU but oversubscribe
GPU memory. In metagenomic sequencing, a hash table is used as the genome
database that is later used for query processing. The size of the entire dataset
can grow to 50+~GB, easily overflowing the GPU capacity. Skip lists are
popularly used in
databases~\cite{rocksdb-meta-2012,leveldb-google-2011,pebblesdb-sosp-2017},
networking~\cite{deterministic-skiplists-soda-1992}, persistent
memory~\cite{listdb-osdi-2022,phast-tpds-2022}, concurrency
control~\cite{kiwi-ppopp-2017,nitro-vldbe-2016,jellyfish-middleware-2020}, and
for representing skewed data access
patterns~\cite{biased-skiplist-algorithm-2001,biased-skiplist-algorithmica-2005}.
In this section, we briefly review existing approaches for designing hash tables and skip lists for GPUs without memory oversubscription. We are not aware of existing work that focuses on improving performance for UVM-enabled GPU data structures.

\paragraph*{Hash Table}

Hash tables are popularly used for tracking key-value relationships in a
sparse domain because of their expected constant time complexity for
different operations like insert and search. The performance of a hash
table depends on the memory layout, access pattern dictated by the input
operation string, quality of the hash functions, and other design decisions
like collision resolution and dynamic resizing techniques. Different
approaches~\cite{art-of-multiprocessor-programming-v2,hashtable-survey-cpu} have been explored for CPUs over several decades and
many libraries provide efficient CPU implementations with different features
(e.g., concurrent hash maps from \textsf{oneTBB}\footnote{\url{https://github.com/uxlfoundation/oneTBB}} and \textsf{Folly}\footnote{\url{https://github.com/facebook/folly}}).

Hash table designs for GPUs can be divided into static~\cite{cudpp-hashtable,gpu-stadium-hashing,gpu-parallel-hashing,lockfree-ds-icpads-2012,warpcore-hipc-2020,hashgraph} and dynamic~\cite{slabhash-ipdps-2020,dachash-sbacpad-2021,warpdrive,bght,gph-pamd-2025,dycuckoo} designs. The container array in the
hash table is not resizable in the static design, which implies that the
former category of approaches make a worst-case estimate on the number of
unique keys before allocating memory. Static memory allocation without
knowing an accurate estimate of the number of unique keys can waste memory.

GPU data structure kernels provide bulk APIs that work on multiple data
elements in parallel to avoid frequent expensive kernel launches~\cite{slabhash-ipdps-2020,warpcore-hipc-2020,gfsl-pact-2017,gph-pamd-2025}. A natural division of work in bulk-API- based kernels is to assign an
operation on a (key, value) pair to a CUDA thread. However, given the
potentially irregular input-dependent access pattern of the data structures,
assigning a different key to each thread in a warp leads to branch and
memory divergence. The Warp-Cooperative Work-Sharing (WCWS) strategy,
popularly followed by recent approaches~\cite{slabhash-ipdps-2020,warpcore-hipc-2020,warpdrive,gfsl-pact-2017,gpu-stadium-hashing}, tries to ameliorate
branch and memory divergence by assigning an operation to a warp.

\begin{wrapfigure}{r}{0.5\textwidth}
    \centering
    \includegraphics[width=0.85\linewidth]{"./figs/slabhash.drawio.pdf"}
    \caption{SlabHash layout with super-nodes; {\tt kv} denotes a (key, value) pair}
    \label{fig:slabhash}
\end{wrapfigure}

\emph{SlabHash}~\cite{slabhash-ipdps-2020} uses closed addressing for conflict
resolution. Unlike traditional linked lists, each bucket in SlabHash is
organized as a list of super-nodes where each super-node consists of multiple
key-value pairs including a pointer to the next node
(Figure~\ref{fig:slabhash}). A super-node is sized based on the warp size (e.g.,
32 for NVIDIA GPUs). SlabHash requires a custom memory allocator to handle
frequent and concurrent dynamic allocation of super-nodes to support resizing.
\emph{DACHash}~\cite{dachash-sbacpad-2021} extends SlabHash to sort the input
sequence for better access locality, and dynamically assigns work to either each
thread or a warp based on the number of super-nodes in the relevant bucket.

\emph{WarpCore}~\cite{warpcore-hipc-2020} proposes a static hash table, i.e.,
the table is not dynamically resized, and uses open addressing for resolving
collisions. WarpCore assigns an operation on a (key, value) pair to one
cooperative group, and all threads in a cooperative group participate to carry
out the operation (\eg, performs a parallel probe to resolve conflicts).
WarpCore may potentially suffer from high overheads in case there is a large
proportion of negative queries (e.g., failed delete or search operations) at
high load factors because it has to keep iterating through the hash table till
it encounters an empty slot~\cite{bght}.

DyCuckoo~\cite{dycuckoo} divides the hash table array into $d$ smaller tables of equal capacity and uses Cuckoo hashing~\cite{cuckoo-hashing} to lower the cost of rehashing when the table is resized. The hash table is downsized or upsized only if the load factor crosses respective thresholds.
% Thus, DyCuckoo improves the amortized cost of negative queries and dynamic resizing.

GPH~\cite{gph-pamd-2025} improves performance by using a prefect hashing scheme.
%It uses a cost model to gauge the performance of the
%hash table based on the occupancy, number of instructions executed, and the
%cycles per instruction.
GPH divides the hash table into $n$ contiguous buckets,
and stores the location of each bucket in an auxiliary array in per-SM shared memory.
GPH uses three hash functions to map a key $k$ to a bucket.
%calculate the
%destination bucket for a key $k$.
The first function maps $k$ to an
auxiliary array index $I$. The second function calculates a virtual
bucket ID using $k$ and the mapping stored at index $I$. The third
function computes the destination bucket.
The auxiliary array allows GPH to
perform an operation with a fixed number of probes.
However, a small fraction of insert operations may fail due to collision.

\begin{figure}[hb]
    \centering
    \includegraphics[width=0.7\columnwidth]{"./figs/basic-skiplist-design-simplified.drawio.pdf"}
    \caption{Layout of a traditional skip list. The data layer shows the keys in stored in the skip list. The dotted bold black line shows the sequence of pointers traversed while searching for the key 11.}
    \label{fig:skiplist-design}
\end{figure}

\paragraph*{Skip List}

A skip list~\cite{skiplist-cacm-1990} is a probabilistic data structure that
allows efficient insertion, deletion, and search operations on sorted linked
lists.
% Rebalancing height-balanced trees is challenging while supporting
% concurrent operations, and hence skip lists are used as a simpler alternative to
% a deterministically balanced tree.
A skip list is organized as a hierarchy of
linked lists, where higher-level lists are always contained in lower-level
lists. The list at the lowest level ($L_{0}$) contains the data elements in a
sorted order. We refer to the lowest level $L_{0}$ as the data layer, and the
upper levels as the index layers. Figure~\ref{fig:skiplist-design} shows the
structure of a skip list traditionally used in CPUs. A skip list employs skip
pointers for efficiently searching through the data structure (denoted by dense
black arrows in Figure~\ref{fig:skiplist-design}). Unlike balanced trees, a skip
list does not require expensive rebalancing after mutation operations while
achieving an expected logarithmic search time. Xing \etal survey different
variants of skip lists~\cite{skiplist-survey-2025}.

Misra and Chaudhuri~\cite{lockfree-ds-icpads-2012} ported the linearizable
construction of a lock-free skip list algorithm described by Herlihy and
Shavit~\cite{art-of-multiprocessor-programming-v2} to GPUs. Moscovici
\etal~\cite{gfsl-pact-2017} proposed GPU-friendly skip list (GFSL) which tries
to reduce the memory footprint and performance penalty due to uncoalesced global
memory accesses on a GPU. GFSL also uses the concept of super-nodes for linked
lists at each level. Figure~\ref{fig:gfsl-node} shows the organization of a
super-node in the index and data layers in GFSL. Each super-node in GFSL
contains $N + 2$ (key, value) pairs. A value field in an index layer $i$
(\code{DownPtr}) stores a pointer to the relevant super-node in the layer $i-i$,
and stores the data associated with the corresponding key in the data layer. The
last two pairs are used to store the maximum key in the node (\code{MaxKey}), a
pointer to the next node at the same level (\code{NextPtr}), and a lock to
ensure mutual exclusion during concurrent mutations.
Figure~\ref{fig:gfsl-design} shows a schematic of the GFSL design where each
super-node accommodates two keys. The blue dashed arrow shows the path taken
while searching for the key 11. A super-node is split if an insert operation
finds the node to be full, and a new super-node is added to the skip list
containing half of the values from the split parent super-node. In contrast to
traditional skip list, a key in GFSL is raised to the higher level only on a
node split operation. Like hash tables, the super-node is sized to match the
size of a warp (\ie, $N + 2$ is 32).

\begin{figure}[t]
    \centering
    \includegraphics[width=0.9\columnwidth]{"./figs/gfsl-node-overview1.drawio.pdf"}
    \caption{Super-nodes in index and data layers in GFSL}
    \label{fig:gfsl-node}
\end{figure}

\begin{figure}[b]
    \centering
    \includegraphics[width=0.7\columnwidth]{"./figs/gfsl-design-overview-sb.drawio.pdf"}
    \caption{GFSL with two keys per super-node
        % and search of key 11
    }
    \label{fig:gfsl-design}
\end{figure}

\section*{Motivation}

We discuss UVM-based implementations of hash table and skip list that we use as
the baseline for this work as a first step toward enabling memory
oversubscription.

\paragraph*{Hash Table with UVM Support}

Unlike WarpCore which uses open-addressing for resolving collisions,
closed-chaining-based designs (\eg, SlabHash~\cite{slabhash-ipdps-2020}) are
sensitive to the memory allocator and requires nontrivial extensions. On the
other hand, probing the hash table array with open-addressing has potentially
better data locality compared to closed-chaining~\cite{warpcore-hipc-2020}.
Hence, we have decided to use WarpCore~\cite{warpcore-hipc-2020} as the baseline
to support the UVM programming model. We will refer to our
UVM-enabled baseline hash table design as \htuvm.

\Htuvm's hash table is static, similar to WarpCore. A limitation of a static
design is that it requires an estimate of the maximum number of insert
operations to avoid expensive resizing at run time. Therefore, \htuvm conservatively
allocates an array of size \code{$\lceil \mathtt{(NUM\_INS/LF)} \rceil$}
in global memory using managed memory allocation of CUDA, where
\code{NUM\_INS} is the number of insert operations to be performed on the
hash table and \code{LF} indicates the fraction of filled entries.
\Htuvm implements three kernels: \code{insert()}, \code{delete()}, and
\code{search()}.
% In addition to a pointer to the hash table in GPU global memory
% and an array of input keys to be operated on, each kernel takes additional
% arguments. The \code{insert()} kernel also takes as input an auxiliary array of
% values corresponding to the input keys that are to be inserted, while the
% \code{delete()} and \code{search()} kernels pass an output array as parameter
% which tracks the values (corresponding to the input keys) that were operated on.
The operations in \htuvm can be invoked both for a single key or for a batch of
keys; the latter will avoid the overhead of repeated kernel launches. \Htuvm
assigns each operation to a cooperative group (CG) which is commonly referred to
as warp-cooperative work scheduling (WCWS), and has been used in prior work~\cite{gfsl-pact-2017,warpcore-hipc-2020}. In addition, we have incorporated managed memory allocation to enable a vanilla UVM support.

\paragraph*{Skip Lists with UVM Support}

We extend the GFSL algorithm~\cite{gfsl-pact-2017} with UVM support to enable
oversubscribing GPU memory, and refer to this design as \skipuvm. The
implementation supports \code{insert()}, \code{delete()},
\code{findPredecessor()}, and \code{search()} operations, similar to existing
skip list implementations.
To enable a vanilla UVM support, we incorporated managed memory allocation for
the super-nodes in GFSL. 

\paragraph*{Optimizing Data Movement}

The vanilla implementation performs very poorly with more than 50\%
oversubscription, because of the random memory access pattern in these data
structures. A general good-quality hash function in a hash table distributes the
memory accesses across different buckets, which leads to touching random pages
depending on the input key pattern. Oversubscription further leads to many pages
storing hash table data to get evicted (based on variants of least-recently-faulted replacement
policy implemented by the contemporary NVIDIA GPUs). This gives rise to frequent
far faults. Far faults are expensive because they incur overheads of page
unmapping on the CPU, page migration to the GPU, and installing page table
entries in the GPU. Multiple concurrent warps (or CGs) stall waiting for their
respective far faults to be resolved, thereby hurting
performance~\cite{irregular-page-walks-isca-2018,neighborhood-aware-addr-translation-micro-2018}.
For example, in our experiments with the baseline UVM hash table, the \code{insert()} kernel does not
finish within even 24 hours with 100\% oversubscription on a GPU with 16 GB
memory.

We optimize data movement
% in the vanilla UVM implementation
by incorporating
% different combinations of
memory advise and
% asynchronous
prefetching hints~\cite{cuda-programming-guide}.
% The best hint combination found through empirical evaluation turns out to be
% \code{cudaMemAdviseSetAccessedBy} for the hash table array, and both
% \code{cudaMemAdviseSetAccessedBy} and \code{cudaMemPrefetchAsync} for the input
% array containing keys to be operated on as well as the auxiliary array of values
% used by the \code{insert()} kernel.
The best combination of hints for the \code{insert()} kernel found through empirical evaluation is
\code{cudaMemAdviseSetAccessedBy} and \code{cudaMemPrefetchAsync} for the hash table array and the input array.
The first hint indicates that the hash table array, the input array, and the auxiliary array will be predominantly accessed by the GPU, allowing the runtime
to pre-map the page table entries on the GPU without actually migrating the data. The \code{cudaMemPrefetchAsync} hint
allows the runtime to asynchronously prefetch the hinted pages before the kernel starts,
thereby reducing the initial far faults.
%An empirical evaluation shows that the
%performance with these hints improves performance compared to the baseline UVM
%implementation.
We also experimented with the \code{cudaMemAdviseSetPreferredLocation} and \code{cudaMemAdviseSetReadMostly} hints.
%, but did not find them useful for improving
% the hash table performance.
The \code{cudaMemAdviseSetPreferredLocation} hint indicates a preferred location (CPU or GPU) for a memory region, but it does not improve performance
because the hash table pages anyway get evicted due to oversubscription.
% and then be migrated back on future access.
The \code{cudaMemAdviseSetReadMostly} hint is not applicable as we update the hash table across kernel launches.

For skip lists, we similarly experiment with the memory advise and prefetch
hints, and find that only prefetch hints are useful for the super-nodes of
SL-UVM and the input key array. We do not use any hints for the auxiliary value
array.

\paragraph*{Shortcomings of the Vanilla UVM Design}

% We now point out the bottlenecks faced by these UVM-based data structure designs
% and why existing techniques that aim to optimize UVM performance fail to address
% them.

% We further study the scalability of a real-world application that uses \htuvm and \skipuvm to implement key-value stores.
We use our hash table and skip list designs, optimized with NVIDIA-provided data migration hints, to implement key-value stores
in a real-world application, \emph{metagenomics classification}~\cite{cuclark,metacache-icpp-2021}, and study the scalability. % Both applications
% rely heavily on key-value stores to efficiently process large-scale genomics
% datasets.
% The metagenomics classifier divides a genome into fixed substrings of length $k$~(called a \textit{k-mer}) and encodes each nucleotide~(Adenine, Thymine, Guanine, Cytosine) of the substring using two bits to form the key $K$ representing the substring.
% Each key $K$ is associated with a \textit{taxon\_id}, which uniquely identifies an organism and acts as the value of the key.
% The classifier consists of two phases, namely, construction of the key-value pair database involving insert operations and the classification phase
% involving search operations.
% The k-mer counting application tracks the frequency of each k-mer in a genome.
% Instead of the \textit{taxon\_id}, the frequency of a k-mer serves as its value.
In our evaluation, we use genome sequences from the NCBI dataset~\cite{refseq-ncbi} for building and querying the database.

Figure~\ref{fig:App-uvm-perf-motivation} shows the scaling of a metagenomics classifier~\cite{metacache-icpp-2021} with data placement and prefetching hints enabled.
%As shown in Figures~\ref{fig:App-SL-motivation} and~\ref{fig:App-HT-motivation},
Both \htuvm and \skipuvm exhibit \emph{poor} scaling behavior in the two key phases of the metagenomics classifier, namely, \textsf{Build}~(uses the \code{insert()} kernel) and \textsf{Classify}~(uses the \code{search()} kernel). As the degree of memory oversubscription increases from 0\% to higher levels, the run time of both the phases increases non-linearly. % The results shown correspond to executions of \htuvm and \skipuvm with memory hints enabled (Figure~\ref{fig:ht-uvm-perf-motivation-w-hints});
Without data placement and prefetching hints, the application run time degrades dramatically, reaching the order of hours.

\begin{figure}[h!]
  \centering
    \begin{subfigure}[b]{0.48\columnwidth}
      \centering
      \includegraphics[width=0.7\linewidth]{"./figs/Applications_Motivation_HT.pdf"}
      \caption{Scaling with \htuvm}
      \label{fig:App-HT-motivation}
    \end{subfigure}%
    \hfill%
    \begin{subfigure}[b]{0.48\columnwidth}
      \centering
      \includegraphics[width=0.7\linewidth]{"./figs/Applications_Motivation_SL.pdf"}
      \caption{Scaling with \skipuvm}
      \label{fig:App-SL-motivation}
    \end{subfigure}
    \caption{Scaling of a metagenomics classifier application~\cite{metacache-icpp-2021} that uses \htuvm or \skipuvm with memory oversubscription}
    \label{fig:App-uvm-perf-motivation}
\end{figure}


These results demonstrate that far page faults and page migration significantly
impact the performance of real-world applications under memory oversubscription,
highlighting their importance in everyday use cases. Existing techniques, such
as prefetching, page placement hints, and page replacement policies, optimize memory usage, but often
fail to be effective when applications with irregular access patterns
oversubscribe GPU
memory~\cite{uvmsmart-isca-2019,hpe-tcad-2020,cppe-ipdps-2020,fine-grain-quantitative-uvm-taco24,in-depth-analysis-uvm-sc21}.
%  Prefetching mitigates
% these issues by proactively loading future-accessed pages, with NVIDIA employing
% a tree-based neighborhood prefetcher~\cite{uvmsmart-isca-2019}. Prior work has
% explored page eviction policies to reduce page faults under memory
% oversubscription, such as hierarchical page eviction and reserved
% LRU~\cite{uvmsmart-isca-2019,hpe-tcad-2020,cppe-ipdps-2020}.
Applications with irregular access patterns, \eg, graph-based applications and
concurrent GPU data structures that require pointer chasing, have poor locality,
making prefetching counterproductive.
Thus, it is imperative to come up with targeted novel data structure designs that can scale with the input size.

\section*{Problem Statement and Proposed Work}

We propose to work on designing efficient GPU data structures to handle use
cases where applications oversubscribe GPU memory. In particular, we will focus
on the following data structures popularly used in GPU kernels: hash table, skip
list, approximate membership query filter, and B- and B\textsuperscript{+}-tree.
These data structures are characterized by random accesses, and hence encounter
frequent pages faults and unnecessary data migrations between the CPU and the
GPU with the UVM programming model.

In the following, we outline our proposal's key research objectives which aim
to explore efficient GPU data structure designs for oversubscription.

\begin{itemize}
    \item Propose a high-performance adaptive AMQ filter design for GPUs,
  \item Extend existing in-memory designs of AMQ and B-tree variants to work with oversubscription (\eg, allocate managed memory with \code{cudaMallocManaged()} and avoid explicit memory copies),
    \item Show limitations in using \nvidia-provided APIs (\eg,  \code{cudaMemAdvise()}) for managing data movement,
    \item Propose more efficient and scalable designs for AMQ and B-tree variants,
    \item Explore performance bottlenecks of UVM applications in a multi-GPU setup.
\end{itemize}

We expand each objective in more detail in the following.

\subsection*{Adaptive Filters for GPUs}

Approximate membership query (AMQ) filters are data structures that maintain a
space-optimized approximate representation of a set or a multiset and allow
efficient probabilistic set membership tests with minimal memory
overhead~\cite{gpu-filters-ppopp23,aqf-sigmod}. AMQ filters significantly reduce
memory overhead costs compared to exact lookup mechanisms, making them
well-suited for large-scale applications. AMQ filters are popularly used in
databases, networks, and file systems, where there is a substantial benefit by
quickly filtering negative queries. AMQs though allow queries to occasionally
return a false positive.

A key limitation of static AMQ structures is their vulnerability to adversarial query patterns. An adversary can repeatedly query elements that consistently yield false positives, degrading the reliability of the system. Unlike filters with a predetermined false positive rate, adaptive filters dynamically adjust to evolving workloads and data distributions, thereby improving their accuracy over time.
Adaptive filters modify their internal state in response to negative queries, effectively reducing repeated false positives. Furthermore, adaptive filters demonstrate robust performance under skewed query distributions, which are common in real-world workloads. These properties make them an increasingly important data structure in contemporary large-scale systems. An implicit assumption is that oracle knowledge is available to decide which queries return false positives so that corrective action can be taken.

Among the existing designs, the adaptive quotient filter (AQF)~\cite{aqf-sigmod}, built upon the RSQF quotient filter~\cite{rsqf}, represents the current state-of-the-art to eliminate recurring false positives. On a query, the AQF performs a secondary check with an exact map (possibly stored on disk) to identify a false positive, and then extends the fingerprint remainder by extra bits to reduce the chances of hash collision. 

However, to the best of our knowledge, there has been no prior work on designing adaptive filters specifically optimized for GPU architectures. Although a quotient filter enjoys better locality than a Bloom filter, the adaptive and counting mechanisms in AQF are difficult to parallelize efficiently. A \naive GPU port will likely perform poorly compared to a CPU implementation because of branch and memory divergence. In this proposal, we aim to develop a GPU-based adaptive filter that leverages the parallel processing capabilities of modern GPUs to significantly improve query throughput and reduce latency. Such a design has the potential to substantially accelerate approximate membership queries in high-performance and large-scale environments, thereby broadening the applicability of adaptive filters in data-intensive systems.

\subsection*{Designing Oversubscribed Data Structures}

We have worked on designing alternate hierarchical memory layouts for hash table
and skip list data structures to maximize access locality and hence improve
performance. The insight is that the outer-level containers enable efficient
jumps to desired regions of the data structure, while the inner containers allow
operating on the data. Our preliminary results show that our custom data
structure designs have the potential to substantially improve performance over
the optimized UVM baselines while supporting high degree of GPU memory
oversubscription for many inputs. Our work has been published in OOPSLA 2026~\cite{gpu-ovs-ds-oopsla-2026}.

We want to continue our work on GPU data strutures by exploring better designs for B-tree and its variants and AMQ filters with UVM support, such that they scale with input size.
Scaling concurrent GPU data structures to high degrees of GPU memory
oversubscription requires novel designs that take into consideration the
potential lack of locality in the access pattern. The key insights for
minimizing page faults are to (i) group nearby keys in the input sequence to
enhance the locality of the input, and (ii) translate the locality, if any,
present in the input sequence to the hash table architecture. The first
requirement ensures that keys within the same page are processed together (i.e.,
near in time) leveraging spatial locality. The second requirement implies that
the nearby keys map to nearby memory pages of the data structure. Taken together, the aforementioned two
requirements optimize memory access pattern and facilitate efficient
prefetching, which, in turn, reduces the frequency of on-demand page faults and
subsequent stalls during operations improving overall performance.

Existing techniques, such as prefetching and page replacement policies, optimize
memory usage, but often fail to be effective when applications with irregular
access patterns oversubscribe GPU
memory~\cite{uvm-forest-isca25,uvmsmart-isca-2019,cppe-ipdps-2020,hpe-tcad-2020}.
Applications with irregular access patterns, e.g., graph-based applications and
concurrent GPU data structures that require pointer chasing, have poor locality,
making prefetching counterproductive. Thus, it is imperative to come up with
targeted novel data structure designs that can scale with the input size.

An alternate strategy to deal with oversubscription is to partition the huge
input workload in multiple batches such that each batch fits within the GPU
memory, and then launch multiple kernels. Batching as an optimization is
commonly performed across application domains (\eg, deep learning training).
\emph{WarpDrive}~\cite{warpdrive} proposed a multi-GPU hash table that
partitions large inputs to fit on a GPU. However, such partitioning is
programmer-driven, tedious, error-prone, brittle (\ie, will break for even
larger inputs), and expensive. Furthermore, certain applications may require
specialized GPU computing algorithms that may preclude manual data
manipulation~\cite{perks-isc-2023}.

\subsection*{Exploring Performance on a Multi-GPU Setup}

% The use of GPUs is ubiquitous today for executing massively parallel AI and ML applications that rely on mathematical operations such as SpMV, SpMM, and convolution kernels. 
% The introduction of UVM by Nvidia~\cite{fine-grain-quantitative-uvm-taco24} simplifies memory management by providing a unified address space accessible across devices. 
% Unified addressing also reduces the complexity of explicit memory transfers and enables developers to design scalable applications.

Modern compute nodes, especially in HPC clusters and cloud environments, often include multiple GPUs per node, making efficient data sharing critical for the seamless operation of applications across devices.
The performance of these applications scales well when the access pattern is regular, but suffers severe degradation under irregular access patterns~\cite {gpu-ovs-ds-oopsla-2026,glift-tog-2006}.
While multi-GPU systems can offer better speedup by exploiting the increased aggregate memory capacity of the multiple devices for larger data sets and improved concurrency, the performance is often limited by the communication overhead associated with the data sharing~\cite{gpusm-gpgpuw-2015,multi-gpu-spmv-taco-2024,multigpu-comm-js-2018}.
Recent research highlights both the opportunities and challenges in this area. For example, studies show that UVM-based implementations can outperform manual memory management by up to 1.8× in some multi-GPU workloads, while significantly improving programmability~\cite{multi-gpu-nas-2021}.
At the same time, newer work on multi-GPU data structures—such as sparse matrix frameworks—focuses on fine-grained data partitioning and load balancing to efficiently distribute irregular workloads across GPUs~\cite{multi-gpu-spmv-taco-2024}.
Other research explores alternative models, such as shared-memory multi-GPU architectures, partitioned global address spaces, and object affinity in the application to GPUs~\cite{oasis-hpca-2025}, to reduce communication overhead and improve scalability.
Overall, studying UVM data structures alongside these emerging approaches is essential for building high-performance, scalable applications in domains such as machine learning, scientific computing, and large-scale graph processing.

\section*{Ongoing Work}

Our work \textit{Designing GPU Data Structures for Efficient Memory Oversubscription} has been accepted to appear in OOPSLA 2026~\cite{gpu-ovs-ds-oopsla-2026}. We are currently studying the design of different AMQ proposals with a view to extending to GPUs. Parallely, we are also extending our hash table and skip list work that targeted nodes with a single GPU to a multi-GPU setup.

\clearpage

{\scriptsize
    \bibliographystyle{plain}
    \bibliography{./../references/venue-full.bib,./../references/references.bib}
}

\end{document}
